\documentclass[11pt, a4paper]{article}

% ═══════════════════════════════════════════════════════════
% PACKAGES
% ═══════════════════════════════════════════════════════════
\usepackage[utf8]{inputenc}
\usepackage[T1]{fontenc}
\usepackage[english]{babel}
\usepackage{amsmath, amssymb, amsthm}
\usepackage{mathtools}
\usepackage{geometry}
\usepackage{hyperref}
\usepackage{booktabs}
\usepackage{array}
\usepackage{xcolor}
\usepackage{tcolorbox}
\usepackage{enumitem}
\usepackage{fancyhdr}
\usepackage{titlesec}
\usepackage{caption}
\usepackage{float}

\geometry{margin=2.5cm}
\hypersetup{
    colorlinks=true,
    linkcolor=blue!70!black,
    citecolor=red!70!black,
    urlcolor=blue!60!black
}

% ═══════════════════════════════════════════════════════════
% THEOREM ENVIRONMENTS
% ═══════════════════════════════════════════════════════════
\theoremstyle{definition}
\newtheorem{definition}{Definition}[section]
\newtheorem{proposition}{Proposition}[section]
\newtheorem{remark}{Remark}[section]

% ═══════════════════════════════════════════════════════════
% CUSTOM COMMANDS
% ═══════════════════════════════════════════════════════════
\newcommand{\feature}[1]{\texttt{#1}}
\newcommand{\clip}[1]{\text{clip}_{[0,1]}\!\left(#1\right)}

\tcbuselibrary{skins, breakable}
\newtcolorbox{featurebox}[1]{
    colback=blue!3, colframe=blue!50!black, fonttitle=\bfseries,
    title={#1}, breakable, arc=2mm
}
\newtcolorbox{refbox}{
    colback=gray!5, colframe=gray!50, fonttitle=\bfseries\small,
    title={References \& Justification}, breakable, arc=1mm
}

% ═══════════════════════════════════════════════════════════
% HEADER
% ═══════════════════════════════════════════════════════════
\pagestyle{fancy}
\fancyhf{}
\fancyhead[L]{\small RL\_phase2 -- Feature Engineering}
\fancyhead[R]{\small \thepage}
\renewcommand{\headrulewidth}{0.4pt}

% ═══════════════════════════════════════════════════════════
% TITLE
% ═══════════════════════════════════════════════════════════
\title{
    \vspace{-1cm}
    \textbf{Game-Theoretic Feature Engineering}\\[0.3cm]
    \large Mathematical Foundations for GTO-Based\\
    No-Limit Hold'em Feature Extraction\\[0.3cm]
    \normalsize RL\_phase2 --- 99-Dimensional State Representation
}
\author{RL\_phase2 Project Documentation}
\date{February 2026}

\begin{document}
\maketitle

\begin{abstract}
This document provides a formal mathematical treatment of the 32 game-theoretic features
used in the RL\_phase2 poker agent's state representation. Each feature is derived from
established concepts in Game Theory Optimal (GTO) poker strategy, decision theory, and
combinatorial analysis. We present the mathematical formulation, implementation details,
normalization scheme, and relevant academic references for each feature. The full feature
vector comprises 99 dimensions, of which 67 encode card/position/stack information and
32 encode game-theoretic quantities described herein.
\end{abstract}

\tableofcontents
\newpage

% ═══════════════════════════════════════════════════════════
% SECTION 1: NOTATION
% ═══════════════════════════════════════════════════════════
\section{Notation and Conventions}

\begin{table}[H]
\centering
\begin{tabular}{cl}
\toprule
\textbf{Symbol} & \textbf{Description} \\
\midrule
$P$ & Current pot size (in chips) \\
$B$ & Bet or raise size facing the player \\
$S$ & Player's remaining stack \\
$C$ & Amount to call ($= B$ when facing a bet) \\
$n$ & Number of remaining streets (preflop=3, flop=2, turn=1, river=0) \\
$N$ & Number of active players \\
$\mathcal{E}(h, b)$ & Estimated equity of hand $h$ on board $b$ \\
$\sigma(x)$ & Sigmoid-like normalization: $\sigma_k(x) = \frac{x}{x + k}$ \\
$\text{clip}_{[a,b]}(x)$ & $\max(a, \min(b, x))$ \\
IP & In Position (BTN, CO) \\
OOP & Out of Position (SB, BB, UTG, MP) \\
\bottomrule
\end{tabular}
\caption{Notation used throughout this document.}
\end{table}

All features are normalized to $[0, 1]$ to facilitate neural network training. We use two primary normalization strategies:
\begin{enumerate}
    \item \textbf{Linear rescaling}: $f = \frac{x - x_{\min}}{x_{\max} - x_{\min}}$, when the range is bounded.
    \item \textbf{Sigmoid normalization}: $f = \frac{x}{x + k}$ for $x \geq 0$, which maps $[0, \infty) \to [0, 1)$ with $f(k) = 0.5$.
\end{enumerate}

% ═══════════════════════════════════════════════════════════
% SECTION 2: EV ESTIMATIONS
% ═══════════════════════════════════════════════════════════
\section{Expected Value Estimations (Features 0--2)}

\subsection{EV Call (Feature 0)}

\begin{featurebox}{EV Call --- Expected Value of Calling}
\begin{definition}
The expected value of calling a bet of size $C$ into a pot of size $P$ with equity $\mathcal{E}$ is:
\begin{equation}
    \text{EV}_{\text{call}}^{\text{raw}} = \mathcal{E} \cdot (P + C) - (1 - \mathcal{E}) \cdot C
\end{equation}
Normalized to $[0, 1]$ with $0.5$ representing breakeven:
\begin{equation}
    \feature{ev\_call} = \clip{\frac{\text{EV}_{\text{call}}^{\text{raw}} / M + 1}{2}}
\end{equation}
where $M = \max(P + C,\; C,\; 1)$ is a scaling factor.
\end{definition}
\end{featurebox}

\begin{refbox}
This is the fundamental EV equation from decision theory applied to poker. The formulation follows directly from the definition of expected value as $\text{EV} = \sum_i p_i \cdot v_i$, weighted by the probability of winning (equity) and the payoff in each outcome. See \textbf{Sklansky} \cite{sklansky1999} Ch.~2 for the standard treatment, and \textbf{Chen \& Ankenman} \cite{chen2006} Ch.~3 for the formal game-theoretic framing.
\end{refbox}

\subsection{EV Fold (Feature 1)}

\begin{featurebox}{EV Fold --- Sunk Cost Indicator}
\begin{definition}
The fold cost measures the fraction of the player's total investment already committed to the pot:
\begin{equation}
    I = \max(0,\; P - S)
\end{equation}
\begin{equation}
    \feature{ev\_fold} = \frac{I}{\max(S + I,\; 1)}
\end{equation}
where $I$ is the estimated invested amount.
When $\feature{ev\_fold} = 0$, folding is free (e.g., checking in the big blind). When $\feature{ev\_fold} \to 1$, the player has invested nearly their entire stack, making folding extremely costly.
\end{definition}
\end{featurebox}

\begin{refbox}
While sunk cost should not theoretically affect optimal play (the ``sunk cost fallacy''), the commitment level is a critical input for determining pot odds and stack-to-pot ratio decisions. \textbf{Harrington} \cite{harrington2004} extensively discusses the concept of ``pot commitment'' in Vol.~1, Ch.~7.
\end{refbox}

\subsection{EV Raise (Feature 2)}

\begin{featurebox}{EV Raise --- Expected Value of Raising}
\begin{definition}
The EV of raising incorporates fold equity $\phi$ (the probability that the opponent folds):
\begin{equation}
    \text{EV}_{\text{raise}} = \phi \cdot \underbrace{\frac{P}{P + S}}_{\text{gain if fold}} + (1 - \phi) \cdot \underbrace{\frac{\mathcal{E} \cdot (P + C)}{P + S}}_{\text{gain if called}}
\end{equation}
\begin{equation}
    \feature{ev\_raise} = \clip{\text{EV}_{\text{raise}}}
\end{equation}
\end{definition}
\end{featurebox}

\begin{refbox}
This two-outcome model is a simplification of the full raise EV calculation. The fold equity component comes from \textbf{Sklansky \& Malmuth} \cite{sklansky1999} and is formalized in \textbf{Tipton} \cite{tipton2019}, Ch.~4. The modern treatment in solver analysis (e.g., PioSolver documentation) confirms that raise EV = weighted sum of fold-through equity and showdown equity.
\end{refbox}

% ═══════════════════════════════════════════════════════════
% SECTION 3: FOLD EQUITY & ODDS
% ═══════════════════════════════════════════════════════════
\section{Fold Equity and Odds (Features 3--6)}

\subsection{Fold Equity (Feature 3)}

\begin{featurebox}{Fold Equity}
\begin{definition}
Fold equity $\phi$ is the estimated probability that all opponents fold to an aggressive action. It is computed heuristically based on:
\begin{itemize}
    \item Street (higher on later streets),
    \item Betting pattern aggressiveness,
    \item Relative bet sizing.
\end{itemize}
\begin{equation}
    \phi = f(\text{street}, \text{actions}, B/P)
\end{equation}
\end{definition}
\end{featurebox}

\begin{refbox}
Fold equity is a concept pioneered by \textbf{Sklansky} \cite{sklansky1999} and extensively used in modern GTO solvers. The fact that total EV = showdown equity + fold equity is the fundamental theorem underlying semi-bluff strategy. See \textbf{Janda} \cite{janda2017} Ch.~2 for a modern GTO interpretation.
\end{refbox}

\subsection{Implied Odds (Feature 4)}

\begin{featurebox}{Implied Odds}
\begin{definition}
Implied odds extend pot odds by accounting for expected future winnings on later streets. For a drawing hand with equity $\mathcal{E}$:
\begin{equation}
    \text{IO} = \frac{P + \text{Future Bets}}{C} \approx \frac{P + \lambda \cdot S}{C}
\end{equation}
where $\lambda \in [0, 1]$ is a scaling factor dependent on hand strength and remaining streets. The feature is normalized as:
\begin{equation}
    \feature{implied\_odds} = \clip{\text{IO}_{\text{norm}}}
\end{equation}
Implied odds are high when: (1) the hand is drawing to a strong made hand, (2) the stack is deep, (3) the draw is hidden (e.g., gutshots have better implied odds than flush draws).
\end{definition}
\end{featurebox}

\begin{refbox}
The concept of implied odds originates from \textbf{Sklansky} \cite{sklansky1999}. \textbf{Miller, Sklansky \& Malmuth} \cite{miller2004} provide a thorough treatment in the context of limit hold'em. \textbf{Harrington} \cite{harrington2004} adapts the concept to no-limit play, noting that implied odds increase with effective stack depth and decrease when draws are obvious.
\end{refbox}

\subsection{Reverse Implied Odds (Feature 5)}

\begin{featurebox}{Reverse Implied Odds}
\begin{definition}
Reverse implied odds capture the risk of completing a draw but still losing to a better hand:
\begin{equation}
    \text{RIO} = f(\mathcal{E}, \text{board texture}, \text{hand dominance risk})
\end{equation}
High RIO situations include: second-nut flush draws, low straight draws on monotone boards, and dominated pairs.
\end{definition}
\end{featurebox}

\begin{refbox}
\textbf{Harrington} \cite{harrington2004} Vol.~2 discusses reverse implied odds situations extensively. \textbf{Tipton} \cite{tipton2019} formalizes this as ``post-completion conditional equity loss'' in Ch.~7.
\end{refbox}

\subsection{Commitment Level (Feature 6)}

\begin{featurebox}{Pot Commitment}
\begin{definition}
The commitment level measures how much of the initial stack has been invested:
\begin{equation}
    \feature{commitment} = \clip{\frac{I}{\max(S + I,\; 1)}}
\end{equation}
where $I = \max(0, P - S)$.

When commitment exceeds $\sim 0.33$ (one-third of stack), conventional wisdom dictates that folding becomes suboptimal for most one-pair or better hands.
\end{definition}
\end{featurebox}

\begin{refbox}
The ``point of no return'' in commitment theory is discussed by \textbf{Harrington} \cite{harrington2004} (the ``Harrington M-ratio'') and formalized by \textbf{Chen \& Ankenman} \cite{chen2006}, who show that in a commitment game, a player should commit when $\text{SPR} < \frac{1}{\mathcal{E}}$.
\end{refbox}

% ═══════════════════════════════════════════════════════════
% SECTION 4: RANGE ANALYSIS
% ═══════════════════════════════════════════════════════════
\section{Range Analysis (Features 7--15)}

\subsection{Range Advantage (Features 7--9)}

\begin{featurebox}{Range Advantage Distribution}
\begin{definition}
The range advantage features decompose the player's estimated hand distribution into three categories:
\begin{align}
    \feature{range\_nutted} &= P(\text{hand} \in \text{top 15\% of range}) \\
    \feature{range\_medium} &= P(\text{hand} \in \text{middle 50\% of range}) \\
    \feature{range\_weak} &= P(\text{hand} \in \text{bottom 35\% of range})
\end{align}
These probabilities are estimated from the hand strength $\mathcal{E}(h, b)$ relative to position-dependent range thresholds.
\end{definition}
\end{featurebox}

\begin{refbox}
Range-based thinking is the cornerstone of modern GTO poker. \textbf{Janda} \cite{janda2017} formalizes range decomposition in Ch.~4. \textbf{Pluribus} \cite{brown2019} implicitly encodes range information through its blueprint strategy. The concept of ``nutted vs.\ medium vs.\ weak'' originated with \textbf{Galfond's} range classification system, later adopted by solver-based training tools like GTO Wizard.
\end{refbox}

\subsection{Board Coverage (Features 10--12)}

\begin{featurebox}{Board Coverage}
\begin{definition}
Board coverage measures how well a player's range connects with different regions of the board:
\begin{align}
    \feature{coverage\_high} &= \text{proportion of range that hits high cards} \\
    \feature{coverage\_medium} &= \text{proportion hitting middle cards} \\
    \feature{coverage\_low} &= \text{proportion hitting low cards}
\end{align}
Coverage is estimated based on position preflop ranges and board rank distribution.
\end{definition}
\end{featurebox}

\begin{refbox}
Board coverage analysis is a concept from modern solver-based poker theory. \textbf{Janda} \cite{janda2017} discusses how GTO strategies adjust c-bet frequency based on board coverage in Ch.~5--6. PioSolver documentation explicitly mentions ``range-vs-board interaction'' as a primary driver of flop strategy.
\end{refbox}

\subsection{Polarization (Features 13--15)}

\begin{featurebox}{Range Polarization}
\begin{definition}
Polarization describes the shape of a player's value distribution:
\begin{align}
    \feature{polarized} &\in [0, 1] \quad \text{(strong or weak, little medium)} \\
    \feature{merged} &\in [0, 1] \quad \text{(mostly medium-strength)} \\
    \feature{capped} &\in [0, 1] \quad \text{(no nutted hands in range)}
\end{align}
A polarized range bets big (either nuts or bluffs). A merged range bets small (medium-strength value). A capped range cannot raise credibly.
\end{definition}
\end{featurebox}

\begin{refbox}
The polarized vs.\ merged vs.\ capped trichotomy is central to modern GTO theory. \textbf{Janda} \cite{janda2017} devotes Ch.~3 to this concept: ``Polarized ranges allow geometric sizing, while merged ranges prefer smaller bets.'' \textbf{Tipton} \cite{tipton2019} formalizes this in his ``Linear vs.\ Polar'' betting model (Ch.~5), showing that optimal bet size is directly determined by the shape of the betting range.
\end{refbox}

% ═══════════════════════════════════════════════════════════
% SECTION 5: BETTING PATTERNS
% ═══════════════════════════════════════════════════════════
\section{Betting Pattern Analysis (Features 16--19)}

\begin{featurebox}{Betting Pattern Classification}
\begin{definition}
The betting pattern of the current action sequence is classified into four strategic archetypes:
\begin{align}
    \feature{pattern\_value} &= P(\text{actions consistent with value betting}) \\
    \feature{pattern\_bluff} &= P(\text{actions consistent with bluffing}) \\
    \feature{pattern\_balanced} &= P(\text{actions consistent with balanced play}) \\
    \feature{pattern\_exploitative} &= P(\text{actions consistent with exploitative play})
\end{align}
These are estimated from the action history (bet sizes, check/raise patterns) on the current street.
\end{definition}
\end{featurebox}

\begin{refbox}
Betting pattern recognition is a key component of opponent modeling. \textbf{Billings et al.} \cite{billings2002} introduced action pattern matching in the Loki poker bot. The value/bluff/balanced/exploitative taxonomy aligns with \textbf{Janda's} \cite{janda2017} classification in Ch.~2 and is implicit in solver output analysis (e.g., the bluff-to-value ratio).
\end{refbox}

% ═══════════════════════════════════════════════════════════
% SECTION 6: GTO FUNDAMENTALS
% ═══════════════════════════════════════════════════════════
\section{GTO Fundamentals (Features 20--23)}

\subsection{Minimum Defense Frequency (Feature 20)}

\begin{featurebox}{MDF --- Minimum Defense Frequency}
\begin{definition}
When facing a bet of size $B$ into a pot of size $P$, the Minimum Defense Frequency is:
\begin{equation}
    \boxed{\text{MDF} = \frac{P}{P + B}}
\end{equation}
This represents the minimum fraction of hands a player must continue with (call or raise) to prevent the opponent's pure bluffs from being automatically profitable.
\end{definition}

\begin{proposition}
If a player folds more than $1 - \text{MDF}$ of their range, the opponent achieves positive expected value by bluffing with any two cards.
\end{proposition}
\begin{proof}
The opponent's bluff EV is $\text{EV}_{\text{bluff}} = P \cdot P(\text{fold}) - B \cdot P(\text{call})$. Setting $\text{EV}_{\text{bluff}} = 0$: $P \cdot (1 - \text{MDF}) = B \cdot \text{MDF}$, which yields $\text{MDF} = P / (P + B)$. \qed
\end{proof}

\textbf{Implementation:}
\begin{equation}
    \feature{mdf} = \begin{cases} P / (P + B) & \text{if } B > 0 \\ 0 & \text{if } B = 0 \text{ (not facing a bet)} \end{cases}
\end{equation}
\end{featurebox}

\begin{refbox}
MDF is a foundational result in game theory applied to poker. The formal derivation appears in \textbf{Chen \& Ankenman} \cite{chen2006}, Ch.~12 (``Defense Frequencies''). \textbf{Janda} \cite{janda2017} presents it as ``the most important equation in poker'' (Ch.~2). The concept directly derives from \textbf{von Neumann \& Morgenstern's} \cite{vonneumann1944} minimax theorem: in a zero-sum game, each player must prevent the opponent from having a dominant strategy.

Typical MDF values in practice:
\begin{center}
\begin{tabular}{cc}
\toprule
Bet Size (relative to pot) & MDF \\
\midrule
$\frac{1}{3}P$ & 75\% \\
$\frac{1}{2}P$ & 67\% \\
$\frac{2}{3}P$ & 60\% \\
$1P$ (pot-size) & 50\% \\
$2P$ (overbet) & 33\% \\
\bottomrule
\end{tabular}
\end{center}
\end{refbox}

\subsection{Alpha --- Bluff Breakeven Threshold (Feature 21)}

\begin{featurebox}{Alpha ($\alpha$) --- Bluff Breakeven Frequency}
\begin{definition}
Alpha is the minimum fold frequency required from the opponent for a bluff of size $B_{\text{eff}}$ to break even:
\begin{equation}
    \boxed{\alpha = \frac{B_{\text{eff}}}{P + B_{\text{eff}}}}
\end{equation}
where $B_{\text{eff}}$ is the effective bet size (last aggressive action or amount to call).

A bluff is $+$EV when $\text{fold equity}\ \phi > \alpha$.
\end{definition}

\textbf{Relationship to MDF:}
\begin{equation}
    \alpha = 1 - \text{MDF}
\end{equation}
Alpha and MDF are two sides of the same coin: MDF constrains the defender, while $\alpha$ constrains the aggressor.
\end{featurebox}

\begin{refbox}
The $\alpha$ parameter is derived identically to MDF but from the bluffer's perspective. \textbf{Sklansky} \cite{sklansky1999} introduced the concept of ``bluff break-even point'' in \textit{The Theory of Poker}. \textbf{Chen \& Ankenman} \cite{chen2006} provide the formal game-theoretic proof in their AKQ toy game analysis (Ch.~5--6), showing that in Nash equilibrium, both the bluffing and defending frequencies are uniquely determined by the pot geometry.
\end{refbox}

\subsection{Bet-to-Pot Ratio (Feature 22)}

\begin{featurebox}{Bet-to-Pot Ratio (Normalized)}
\begin{definition}
The bet-to-pot ratio captures the relative sizing of the current bet, normalized using a sigmoid function:
\begin{equation}
    r = \frac{B_{\text{eff}}}{P}
\end{equation}
\begin{equation}
    \feature{bet\_to\_pot} = \frac{r}{1 + r} = \sigma_1(r)
\end{equation}

This maps the ratio to $[0, 1)$:
\begin{center}
\begin{tabular}{cc}
\toprule
Raw Ratio $r$ & Feature Value \\
\midrule
0 (no bet) & 0.00 \\
$\frac{1}{3}$ & 0.25 \\
$\frac{1}{2}$ & 0.33 \\
1 (pot-size) & 0.50 \\
2 (2x pot) & 0.67 \\
$\infty$ (all-in) & $\to 1.00$ \\
\bottomrule
\end{tabular}
\end{center}
\end{definition}
\end{featurebox}

\begin{refbox}
Bet sizing theory is extensively covered in \textbf{Tipton} \cite{tipton2019}, Ch.~5--6. The sigmoid normalization $\sigma_k(x) = x/(x+k)$ is a standard technique for unbounded positive quantities in ML feature engineering \cite{bishop2006}. In GTO poker, bet sizing conveys information about range composition: larger bets signal more polarized ranges \cite{janda2017}.
\end{refbox}

\subsection{Pot Geometry (Feature 23)}

\begin{featurebox}{Pot Geometry --- Multi-Street Sizing Optimality}
\begin{definition}
Geometric bet sizing determines the optimal per-street bet size to deploy the entire stack over $n$ remaining streets. The geometric fraction is:
\begin{equation}
    \boxed{g = \left(\frac{S + P}{P}\right)^{1/n} - 1}
\end{equation}
The optimal geometric bet on the current street is $B_{\text{geo}} = P \cdot g$.

The feature compares the actual bet to the geometric optimal:
\begin{equation}
    \feature{pot\_geometry} = \begin{cases}
        \sigma_1\!\left(\dfrac{B_{\text{eff}}}{B_{\text{geo}}}\right) & \text{if } B_{\text{eff}} > 0 \text{ and } B_{\text{geo}} > 0 \\[6pt]
        0.5 & \text{otherwise (neutral)}
    \end{cases}
\end{equation}

\begin{itemize}
    \item $\feature{pot\_geometry} = 0.5$: sizing exactly matches geometric optimal.
    \item $\feature{pot\_geometry} > 0.5$: over-geometric (polarized sizing).
    \item $\feature{pot\_geometry} < 0.5$: under-geometric (merged sizing).
\end{itemize}
\end{definition}
\end{featurebox}

\begin{refbox}
Geometric bet sizing was first formalized by \textbf{Tipton} \cite{tipton2019} in Ch.~6 (``Multi-Street Games''). The key insight is that in a polarized betting game with $n$ streets remaining, the optimal strategy distributes the stack evenly across streets using the geometric series. This result is confirmed by solver outputs: PioSolver's ``Geometric Mode'' uses exactly this formula. \textbf{Janda} \cite{janda2017} discusses geometric sizing in the context of building a pot for a river shove (Ch.~3).

The derivation proceeds from requiring that the total stack is deployed as:
\begin{equation}
    P \cdot (1 + g)^n = S + P
\end{equation}
Solving for $g$ yields the formula above. For example, with $S = 300$, $P = 100$, and $n = 2$ streets:
\[
    g = \left(\frac{400}{100}\right)^{1/2} - 1 = 2 - 1 = 1 \quad \Rightarrow \text{bet pot on each street.}
\]
\end{refbox}

% ═══════════════════════════════════════════════════════════
% SECTION 7: TIER 1 --- ADVANCED GTO
% ═══════════════════════════════════════════════════════════
\section{Tier 1: Advanced GTO Features (Features 24--26)}

\subsection{Equity Realization (Feature 24)}

\begin{featurebox}{Equity Realization}
\begin{definition}
Equity realization (EQR) is the fraction of raw equity that a hand actually converts into chips, accounting for positional and strategic factors:
\begin{equation}
    \boxed{\text{EQR} = \beta_{\text{pos}} \cdot \delta_{\text{hand}} \cdot \mu_{\text{mw}}}
\end{equation}
where:
\begin{align}
    \beta_{\text{pos}} &= \begin{cases} 0.95 & \text{if IP (BTN, CO)} \\ 0.70 & \text{if OOP (SB, BB, UTG, MP)} \end{cases} \\[6pt]
    \delta_{\text{hand}} &= \begin{cases}
        0.85 & \text{single draw (flush or straight draw)} \\
        0.92 & \text{combo draw (flush + straight)} \\
        1.05 & \text{nutted hand } (\mathcal{E} > 0.8) \\
        0.75 & \text{weak hand } (\mathcal{E} < 0.3) \\
        1.00 & \text{otherwise}
    \end{cases} \\[6pt]
    \mu_{\text{mw}} &= \max\!\left(0.7,\; 1 - 0.05 \cdot (N - 2)\right)
\end{align}
\end{definition}
\end{featurebox}

\begin{refbox}
Equity realization is a concept that emerged from solver analysis in the mid-2010s. \textbf{GTO Wizard} (\url{https://gtowizard.com}) and \textbf{Monker Solver} documentation define EQR as:
\[
    \text{EQR} = \frac{\text{Actual EV}}{\text{Raw Equity} \times \text{Pot}}
\]
The position-based coefficients ($\beta_{\text{IP}} \approx 0.95$, $\beta_{\text{OOP}} \approx 0.70$) are empirically validated by solver simulations. \textbf{Brokos} \cite{brokos2014} discusses EQR in his treatment of preflop play (Ch.~4): ``Being in position allows you to realize nearly all of your equity, while being out of position costs you roughly 20--30\%.'' The multiway penalty is documented in \textbf{Janda} \cite{janda2017}, Ch.~8: equity dilution in multiway pots follows approximately $\text{EQR}_{\text{mw}} \approx \text{EQR}_{\text{HU}} \times (1 - 0.05(N-2))$.
\end{refbox}

\subsection{Blocker Effects (Feature 25)}

\begin{featurebox}{Blocker Effects}
\begin{definition}
Blockers measure how much the player's hole cards reduce the number of strong hand combinations available to opponents. The blocker score is computed as a sum of four components:

\begin{equation}
    \feature{blockers} = \clip{\sum_{i=1}^{4} w_i \cdot \mathbb{1}[\text{condition}_i]}
\end{equation}

\begin{enumerate}
    \item \textbf{Nut flush blocker} ($w_1 = 0.4$): Hero holds $A_s$ where $s$ is the dominant board suit with $\geq 3$ cards. Reduces opponent's nut flush combos from $C$ to $0$.
    
    \item \textbf{High card blocker} ($w_2 = 0.2$): Hero holds the same rank as the board's highest card (for ranks $\geq Q$). Reduces top-pair combinations by 50\% (from $\binom{3}{1} = 3$ combos to $\binom{2}{1} = 2$).
    
    \item \textbf{Set blocker} ($w_3 = 0.1$): Hero holds a card matching a board card (rank $\geq 10$). Reduces opponent's set combos from $\binom{3}{2} = 3$ to $\binom{2}{2} = 1$.
    
    \item \textbf{Straight blocker} ($w_4 = 0.1$): Hero holds cards that fill gaps between board cards, reducing straight combinations.
\end{enumerate}
\end{definition}
\end{featurebox}

\begin{refbox}
Blocker analysis is a central concept in advanced GTO play. \textbf{Janda} \cite{janda2017} devotes Ch.~9 to blockers: ``The most important blocker in poker is the nut flush blocker [...] holding $A\spadesuit$ when the board has three spades immediately removes 100\% of the nut flush combinations from the opponent's range.''

The combinatorial justification: a deck has 4 cards of each rank. If hero holds one, opponents only have $\binom{3}{k}$ instead of $\binom{4}{k}$ combinations. For pairs: $\binom{3}{2} = 3$ vs.\ $\binom{4}{2} = 6$, a 50\% reduction. \textbf{Brokos} \cite{brokos2014} discusses blocker-based bluffing in Ch.~8.

The weights ($w_i$) are calibrated from solver outputs. In GTO Wizard's analysis, nut flush blockers increase bluff profitability by approximately 35--45\%, while set blockers contribute approximately 8--12\%.
\end{refbox}

\subsection{Protection Need (Feature 26)}

\begin{featurebox}{Protection Need}
\begin{definition}
Protection need quantifies the urgency of betting to deny equity to opponents' draws:
\begin{equation}
    \feature{protection} = \begin{cases}
        \displaystyle \left(1 - \frac{|\mathcal{E} - 0.65|}{0.25}\right)^+ \!\cdot\; (1 + D_b) \cdot \tau_s & \text{if } 0.4 \leq \mathcal{E} \leq 0.85 \\[8pt]
        0.1 & \text{if } \mathcal{E} > 0.85 \text{ (nutted)} \\
        0.0 & \text{if } \mathcal{E} < 0.4 \text{ (air)}
    \end{cases}
\end{equation}
where:
\begin{align}
    D_b &= 0.4 \cdot W + 0.3 \cdot K \quad \text{(board danger: $W$ = wet, $K$ = coordinated)} \\
    \tau_s &= \begin{cases} 1.0 & \text{flop} \\ 0.7 & \text{turn} \\ 0.0 & \text{river} \end{cases}
\end{align}
The $(\cdot)^+$ denotes $\max(0, \cdot)$. Peak protection occurs at $\mathcal{E} = 0.65$ (top pair / overpair range).
\end{definition}
\end{featurebox}

\begin{refbox}
Protection betting is one of the key distinctions between GTO and exploitative play. \textbf{Tipton} \cite{tipton2019} formalizes the concept in Ch.~8: ``A bet has protection value when it causes opponents to fold hands that have positive equity against our holding.'' 

The bell-curve shape centered at $\mathcal{E} = 0.65$ reflects the GTO principle that:
\begin{itemize}
    \item Nutted hands ($\mathcal{E} > 0.85$) prefer slow-playing or thin value betting.
    \item Air ($\mathcal{E} < 0.4$) gains nothing from protection (no equity to protect).
    \item Vulnerable strong hands ($\mathcal{E} \approx 0.65$) benefit most from charging draws.
\end{itemize}
The board danger term amplifies protection on wet/coordinated boards, consistent with solver output: c-bet frequency is highest with top pair on wet boards (\textbf{Janda} \cite{janda2017}, Ch.~6). The street multiplier $\tau_s$ captures that protection is most urgent on the flop (two cards to come) and irrelevant on the river (no more draws).
\end{refbox}

% ═══════════════════════════════════════════════════════════
% SECTION 8: TIER 2 --- STRATEGIC
% ═══════════════════════════════════════════════════════════
\section{Tier 2: Strategic Features (Features 27--29)}

\subsection{Nut Advantage (Feature 27)}

\begin{featurebox}{Nut Advantage}
\begin{definition}
Nut advantage measures whether the player's positional range contains a higher density of nutted hands than the opponent's range on the given board texture:
\begin{equation}
    \feature{nut\_adv} = \clip{\nu_{\text{base}} - 0.15 \cdot K - 0.1 \cdot \mathbb{1}[\text{monotone}]}
\end{equation}
where:
\begin{equation}
    \nu_{\text{base}} = \begin{cases}
        0.70 \text{ (IP)} / 0.55 \text{ (OOP)} & \text{Ace-high board} \\
        0.65 \text{ (IP)} / 0.50 \text{ (OOP)} & \text{King-high board} \\
        0.35 \text{ (IP)} / 0.60 \text{ (OOP)} & \text{Low board (max rank} \leq 8\text{)} \\
        0.50 & \text{Mid board}
    \end{cases}
\end{equation}
\end{definition}
\end{featurebox}

\begin{refbox}
Nut advantage is one of the primary drivers of GTO strategy selection, particularly regarding c-bet frequency and sizing. \textbf{Janda} \cite{janda2017}, Ch.~5:

``On Ace-high boards, the preflop aggressor has a significant nut advantage because all the Ax hands are primarily in their range. On low, connected boards, the caller's range has more sets and two-pair combinations, giving them the nut advantage.''

The IP coefficients ($0.70$ on A-high) are derived from the standard 6-max opening/3-betting ranges: BTN opens $\sim$45\% of hands and all Ax suited, while BB defends $\sim$35\% with fewer Ax combos. This is confirmed by \textbf{GTO Wizard's} range-vs-board analysis tools. The reduction for coordinated ($-0.15K$) and monotone ($-0.10$) boards reflects that these textures create more possible combinations, diluting any single player's nut advantage.
\end{refbox}

\subsection{Leverage (Feature 28)}

\begin{featurebox}{Leverage --- Multi-Street Pressure Capacity}
\begin{definition}
Leverage measures the ability to credibly threaten large future bets, as a function of the stack-to-pot ratio (SPR) and remaining streets:
\begin{equation}
    \boxed{\feature{leverage} = \min\!\left(\frac{\text{SPR}}{10},\; 1\right) \cdot \frac{n}{3}}
\end{equation}
where $\text{SPR} = S / P$ and $n$ is the number of remaining streets.

\begin{itemize}
    \item $\feature{leverage} = 1.0$: deep stack, preflop (maximum pressure).
    \item $\feature{leverage} = 0.0$: river or all-in (no future threats).
    \item $\feature{leverage} \approx 0.67$: flop with SPR $\geq 10$.
\end{itemize}
\end{definition}
\end{featurebox}

\begin{refbox}
The concept of leverage in poker relates to the ``threat of future bets'' discussed extensively by \textbf{Harrington} \cite{harrington2004} and formalized by \textbf{Tipton} \cite{tipton2019}, Ch.~9. The key insight: a player with a deep stack and multiple streets remaining can make small bets that carry an implicit threat of escalation. \textbf{Chen \& Ankenman} \cite{chen2006} show that in $n$-street games with SPR $= s$, the optimal bluffing frequency depends on $s^{1/n}$, formally linking stack depth, streets, and bluffing credibility.

The normalization at $\text{SPR} = 10$ is conventional: at $\text{SPR} \geq 10$, a player can comfortably make geometric bets across all remaining streets. Below $\text{SPR} = 4$, stack-off decisions are largely mechanical \cite{harrington2004}.
\end{refbox}

\subsection{Effective SPR (Feature 29)}

\begin{featurebox}{Effective Stack-to-Pot Ratio}
\begin{definition}
The stack-to-pot ratio, normalized using a sigmoid centered at $k = 4$:
\begin{equation}
    \boxed{\feature{effective\_spr} = \frac{\text{SPR}}{\text{SPR} + 4} = \sigma_4(\text{SPR})}
\end{equation}

This sigmoid normalization provides useful resolution across the critical SPR range:
\begin{center}
\begin{tabular}{ccc}
\toprule
SPR & Feature Value & Strategic Implication \\
\midrule
1 & 0.20 & Commit zone (top pair = stack-off) \\
2 & 0.33 & Shallow (overpairs commit) \\
4 & 0.50 & Medium (set mining viable) \\
10 & 0.71 & Deep (wide implied odds) \\
20 & 0.83 & Very deep (speculative hands gain) \\
\bottomrule
\end{tabular}
\end{center}
\end{definition}
\end{featurebox}

\begin{refbox}
SPR is one of the most important concepts in no-limit hold'em post-flop play. \textbf{Miller, Sklansky \& Malmuth} \cite{miller2007} formalized SPR in \textit{Professional No-Limit Hold'em: Volume 1}, dedicating an entire chapter to it. Their key thresholds:
\begin{itemize}
    \item SPR $< 3$: top pair is strong enough to commit.
    \item SPR $3$--$6$: overpairs and top pair with good kicker.
    \item SPR $6$--$13$: two-pair and sets are needed.
    \item SPR $> 13$: set mining and speculative hands gain value.
\end{itemize}
The sigmoid normalization at $k = 4$ places the midpoint at the strategically critical SPR$= 4$ boundary, which separates ``shallow'' from ``deep'' play. This choice is supported by solver analysis: at SPR $= 4$, the GTO strategy transitions from ``bet-fold top pair'' to ``bet-call top pair'' in most spots.
\end{refbox}

% ═══════════════════════════════════════════════════════════
% SECTION 9: TIER 3 --- TACTICAL
% ═══════════════════════════════════════════════════════════
\section{Tier 3: Tactical Features (Features 30--31)}

\subsection{Check-Raise Signal (Feature 30)}

\begin{featurebox}{Check-Raise Signal}
\begin{definition}
The check-raise signal detects favorable check-raise spots based on three conditions:

\begin{equation}
    \feature{cr\_signal} = \begin{cases}
        0.7 \cdot (1 - 0.3W) \cdot \mathbb{1}[C > 0] & \text{if OOP AND } 0.65 \leq \mathcal{E} \leq 0.90 \\[4pt]
        0.5 + 0.15 \cdot \mathbb{1}[\text{combo draw}] & \text{if OOP AND draw with } \mathcal{E} < 0.4 \\[4pt]
        \text{above} \times 0.3 & \text{if no bet to face (}C = 0\text{)} \\
        0 & \text{if IP or preflop}
    \end{cases}
\end{equation}

where $W$ = board wetness and $C$ = amount to call.

The check-raise range composition follows the GTO principle:
\begin{equation}
    \frac{|\text{Value CR}|}{|\text{Value CR}| + |\text{Bluff CR}|} \approx 1 - \alpha_{\text{raise}}
\end{equation}
\end{definition}
\end{featurebox}

\begin{refbox}
GTO check-raising frequencies are extensively analyzed in solver outputs. \textbf{Janda} \cite{janda2017}, Ch.~7: ``The optimal check-raise range on the flop contains strong value hands (sets, two-pair) and semi-bluff draws (combo draws, open-enders) in a ratio determined by the raise sizing.''

Key principles reflected in the feature:
\begin{enumerate}
    \item \textbf{OOP requirement}: Check-raising is only possible OOP (you must check first, then raise).
    \item \textbf{Board texture}: Dry boards are better for check-raising because scare cards are less likely on future streets, making continuation more predictable \cite{janda2017}.
    \item \textbf{Semi-bluff selection}: GTO solvers consistently choose combo draws (flush + straight draw) as the primary check-raise bluff because they have the highest equity when called ($\sim$45\%), making them closest to indifferent between being called and seeing a fold \cite{tipton2019}.
\end{enumerate}
\end{refbox}

\subsection{Equity Denial (Feature 31)}

\begin{featurebox}{Equity Denial}
\begin{definition}
Equity denial quantifies the value gained by forcing opponents to fold draws before they can improve:
\begin{equation}
    \feature{eq\_denial} = \begin{cases}
        \clip{0.8\,\mathcal{E} \cdot (1 + D) \cdot (1 + 0.1(N-2))} & \text{if } 0.45 \leq \mathcal{E} \leq 0.85 \text{ and street} \neq \text{river} \\
        0 & \text{otherwise}
    \end{cases}
\end{equation}
where $D = 0.5W + 0.3K + 0.15 \cdot \mathbb{1}[\exists\, \text{flush draw on board}]$ is the draw density of the board.

The equity denial value is zero on the river (no more cards to come) and for air hands (nothing to protect).
\end{definition}

\begin{proposition}
The EV of a protection bet of size $B$ is:
\begin{equation}
    \text{EV}_{\text{protect}} = \phi \cdot \sum_{i \in \text{draws}} \mathcal{E}_i^{\text{opp}} + (1 - \phi) \cdot (\text{EV}_{\text{value}})
\end{equation}
where $\mathcal{E}_i^{\text{opp}}$ is the equity of opponent draw $i$ (which is denied when they fold), and $\phi$ is the probability of folding.
\end{proposition}
\end{featurebox}

\begin{refbox}
Equity denial is the primary reason GTO strategies bet with medium-strength hands. \textbf{Tipton} \cite{tipton2019}, Ch.~8: ``The value of a bet is not limited to how often we get called by worse. A significant component is the equity we deny to hands that fold but would have improved.''

The board draw density $D$ is computed from observable texture features:
\begin{itemize}
    \item Two cards of the same suit ($+0.15$): flush draw possible with 9 outs ($\sim$35\% by river).
    \item Connected board ($+0.3K$): straight draws (4--8 outs, $\sim$17$-$32\% by river).
    \item Wet board ($+0.5W$): multiple draws overlay.
\end{itemize}

The multiway amplification $(1 + 0.1(N-2))$ reflects that with more opponents, the aggregate equity of all opponent draws increases, making protection more valuable. \textbf{Janda} \cite{janda2017} notes that ``in multiway pots, the GTO strategy shifts towards more protection betting and less slow-playing'' (Ch.~8).
\end{refbox}

% ═══════════════════════════════════════════════════════════
% SECTION 10: SUMMARY
% ═══════════════════════════════════════════════════════════
\section{Summary and Feature Vector Structure}

\begin{table}[H]
\centering
\small
\begin{tabular}{ccllc}
\toprule
\textbf{Idx} & \textbf{Tier} & \textbf{Feature Name} & \textbf{Key Formula} & \textbf{Range} \\
\midrule
0 & Base & \texttt{ev\_call} & $(\mathcal{E} \cdot (P+C) - (1-\mathcal{E}) \cdot C) / M$ & $[0,1]$ \\
1 & Base & \texttt{ev\_fold} & $I / (S + I)$ & $[0,1]$ \\
2 & Base & \texttt{ev\_raise} & $\phi \cdot G_f + (1-\phi) \cdot G_c$ & $[0,1]$ \\
3 & Base & \texttt{fold\_equity} & $\phi = f(\text{street}, B/P)$ & $[0,1]$ \\
4 & Base & \texttt{implied\_odds} & $(P + \lambda S) / C$ (norm.) & $[0,1]$ \\
5 & Base & \texttt{reverse\_implied\_odds} & $f(\mathcal{E}, \text{board})$ & $[0,1]$ \\
6 & Base & \texttt{commitment} & $I / (S + I)$ & $[0,1]$ \\
7--9 & Base & \texttt{range\_*} & $P(\text{nutted}|\text{medium}|\text{weak})$ & $[0,1]^3$ \\
10--12 & Base & \texttt{coverage\_*} & Board interaction & $[0,1]^3$ \\
13--15 & Base & \texttt{polar.*} & Range shape & $[0,1]^3$ \\
16--19 & Base & \texttt{pattern\_*} & Action history & $[0,1]^4$ \\
\midrule
20 & GTO & \texttt{mdf} & $P / (P + B)$ & $[0,1]$ \\
21 & GTO & \texttt{alpha} & $B / (P + B)$ & $[0,1]$ \\
22 & GTO & \texttt{bet\_to\_pot} & $r / (1 + r)$ & $[0,1]$ \\
23 & GTO & \texttt{pot\_geometry} & $(B/B_{\text{geo}}) / (1 + B/B_{\text{geo}})$ & $[0,1]$ \\
\midrule
24 & T1 & \texttt{eq\_realization} & $\beta \cdot \delta \cdot \mu$ & $[0,1]$ \\
25 & T1 & \texttt{blockers} & $\sum w_i \cdot \mathbb{1}_i$ & $[0,1]$ \\
26 & T1 & \texttt{protection} & Bell curve $\times D_b \times \tau$ & $[0,1]$ \\
\midrule
27 & T2 & \texttt{nut\_advantage} & $\nu - 0.15K - 0.1M$ & $[0,1]$ \\
28 & T2 & \texttt{leverage} & $\min(\text{SPR}/10, 1) \cdot n/3$ & $[0,1]$ \\
29 & T2 & \texttt{effective\_spr} & $\text{SPR} / (\text{SPR} + 4)$ & $[0,1]$ \\
\midrule
30 & T3 & \texttt{cr\_signal} & OOP $\times$ strength $\times$ board & $[0,1]$ \\
31 & T3 & \texttt{eq\_denial} & $0.8\mathcal{E} \cdot (1 + D)$ & $[0,1]$ \\
\bottomrule
\end{tabular}
\caption{Complete list of the 32 game-theoretic features (indices 67--98 in the full 99-feature vector).}
\end{table}

% ═══════════════════════════════════════════════════════════
% REFERENCES
% ═══════════════════════════════════════════════════════════
\begin{thebibliography}{99}

\bibitem{vonneumann1944}
J.~von Neumann and O.~Morgenstern,
\textit{Theory of Games and Economic Behavior},
Princeton University Press, 1944.
\newblock Foundational work on game theory. Chapter~5 introduces the minimax theorem.

\bibitem{nash1950}
J.~Nash,
``Equilibrium Points in N-Person Games,''
\textit{Proceedings of the National Academy of Sciences}, vol.~36, no.~1, pp.~48--49, 1950.
\newblock Defines Nash equilibrium, the theoretical basis for GTO play.

\bibitem{sklansky1999}
D.~Sklansky,
\textit{The Theory of Poker},
Two Plus Two Publishing, 1999 (revised edition).
\newblock Introduces fundamental theorem of poker, fold equity, pot odds, implied odds.

\bibitem{harrington2004}
D.~Harrington and B.~Robertie,
\textit{Harrington on Hold'em: Expert Strategy for No-Limit Tournaments},
Vols.~1--3, Two Plus Two Publishing, 2004--2006.
\newblock Pot commitment, M-ratio, strategic zones, and multi-street planning.

\bibitem{miller2004}
E.~Miller, D.~Sklansky, and M.~Malmuth,
\textit{Small Stakes Hold'em: Winning Big with Expert Play},
Two Plus Two Publishing, 2004.
\newblock Detailed treatment of implied odds, reverse implied odds, and drawing decisions.

\bibitem{miller2007}
E.~Miller, M.~Sklansky, and D.~Malmuth,
\textit{Professional No-Limit Hold'em: Volume 1},
Two Plus Two Publishing, 2007.
\newblock Introduces and formalizes the Stack-to-Pot Ratio (SPR) concept with strategic thresholds.

\bibitem{chen2006}
B.~Chen and J.~Ankenman,
\textit{The Mathematics of Poker},
ConJelCo, 2006.
\newblock Formal game-theoretic treatment: EV calculations, defense frequencies, commitment theory, AKQ game, $n$-street bluffing theory.

\bibitem{billings2002}
D.~Billings, A.~Davidson, J.~Schaeffer, and D.~Szafron,
``The Challenge of Poker,''
\textit{Artificial Intelligence}, vol.~134, no.~1--2, pp.~201--240, 2002.
\newblock Opponent modeling, action pattern recognition in Loki poker bot.

\bibitem{brokos2014}
N.~Brokos,
\textit{Play Optimal Poker: Practical Game Theory for Every Poker Player},
2019.
\newblock Modern GTO primer: equity realization, blocker-based bluffing, range construction.

\bibitem{janda2017}
M.~Janda,
\textit{Applications of No-Limit Hold'em},
Two Plus Two Publishing, 2017.
\newblock Comprehensive GTO reference: MDF, $\alpha$, polarization, range advantage, nut advantage, board texture analysis, check-raising, geometric sizing, multiway adjustments.

\bibitem{tipton2019}
W.~Tipton,
\textit{Expert Heads Up No Limit Hold'em},
Vols.~1--2, Two Plus Two Publishing, 2019.
\newblock Mathematical treatments of geometric sizing (Ch.~6), polarized vs.\ linear betting (Ch.~5), protection value (Ch.~8), multi-street leverage (Ch.~9).

\bibitem{brown2019}
N.~Brown and T.~Sandholm,
``Superhuman AI for Multiplayer Poker,''
\textit{Science}, vol.~365, no.~6456, pp.~885--890, 2019.
\newblock Pluribus agent: blueprint strategy, Monte Carlo CFR, implicit range encoding.

\bibitem{bishop2006}
C.~Bishop,
\textit{Pattern Recognition and Machine Learning},
Springer, 2006.
\newblock Feature normalization and preprocessing for neural network inputs (Ch.~5).

\end{thebibliography}

\end{document}
